\section{Optimal estimation and interval length}

Identification in \cref{obj:prop:identification} makes the population causal mean a well-defined target of the observed experiment. We now characterize the precision attainable from that experiment as sample size and measurement precision vary. The first criterion is the minimax expected absolute error \leanref{sym:Rn}{\ensuremath{R_n(\sigma)}}, evaluated over the structural class in \cref{obj:def:model-class}.

\begin{definitionv}{def:minimax-risk}[Minimax absolute-error risk $R_{K,n}(\sigma)$]
The minimax expected absolute error is
\[
R_{K,n}(\sigma)
=
\inf_T
\sup_{P\in\mathcal P_{K,\beta,\kappa,\sigma}}
E_{Q_P^n}\!\left[|T-\theta(P)|\right].
\]
Here \(\theta(P)\) is the population causal mean at the prespecified latent dose, and \(Q_P^n\) is the observed-sample law induced by the structural law \(P\), with potential-outcome schedules on the public path space of \cref{obj:def:potential-path-space}. The infimum ranges over measurable, possibly randomized estimators \(T\).
\end{definitionv}

For uncertainty quantification, fix a noncoverage probability \leanref{sym:alpha}{\ensuremath{\alpha\in(0,1/2)}}, corresponding to coverage of at least \(1-\alpha\) under every law in the structural class. The second criterion, \leanref{sym:Hn}{\ensuremath{H_n(\sigma)}}, measures the smallest worst-law expected length among connected intervals satisfying this uniform coverage requirement. The separate finite-sample upper guarantee in \cref{obj:thm:finite-honesty} holds for the wider range \(\alpha\in(0,1)\); the restriction \(\alpha<1/2\) is used by our two-point argument for the honest-length lower bound, which requires \(1-2\alpha\) to be positive; we do not claim that the frontier fails when \(\alpha\ge1/2\).

\begin{definitionv}{def:honest-length}[Honest interval length $H_{K,\alpha,n}(\sigma)$]
The minimum worst-law expected length at fixed noncoverage probability \(\alpha\) is
\[
H_{K,\alpha,n}(\sigma)
=
\inf_{\substack{
C:\ \inf_{P\in\mathcal P_{K,\beta,\kappa,\sigma}}
Q_P^n\{\theta(P)\in C\}\ge 1-\alpha
}}
\sup_{P\in\mathcal P_{K,\beta,\kappa,\sigma}}
E_{Q_P^n}\!\left[\operatorname{length}(C)\right].
\]
The infimum ranges over measurable connected intervals \(C\). Here \(\theta(P)\) is the population causal mean at the prespecified latent dose, and \(Q_P^n\) is the observed-sample law induced by the structural law \(P\), with potential-outcome schedules on the public path space of \cref{obj:def:potential-path-space}.
\end{definitionv}

Both criteria are governed by a common rate \leanref{sym:rho}{\ensuremath{\rho_{n,\sigma}}}, obtained by raising the localization scale \leanref{sym:bns}{\ensuremath{b_{n,\sigma}}} to the mean-smoothness exponent. The effective exponent \leanref{sym:d}{\ensuremath{d}} combines smoothness and local design decay, while the logarithmic sample-size quantity \leanref{sym:Ln}{\ensuremath{L_n}} determines the second measurement-error threshold. Their definitions distinguish three ranges of measurement precision.

\begin{definitionv}{def:frontier-rate}[Frontier rate \(\rho_{n,\sigma}\)]
The frontier rate and its localization scale are
\[
\rho_{n,\sigma}=b_{n,\sigma}^{\beta},
\qquad
b_{n,\sigma}=
\begin{cases}
n^{-1/d},&0\le\sigma\le n^{-1/d},\\[2pt]
\dfrac{\sigma}{\sqrt{\log(e+n\sigma^d)}},
&n^{-1/d}<\sigma\le L_n^{-1/2},\\[6pt]
\dfrac{\log(e+\sigma^2L_n)}{L_n},
&L_n^{-1/2}<\sigma\le1/4.
\end{cases}
\]
Here the effective exponent \(d\) and logarithmic sample-size quantity \(L_n\) are
\[
d=2\beta+\kappa+1,
\qquad
L_n=\log(en).
\]
\end{definitionv}

The main characterization, \cref{obj:thm:uniform-frontier}, gives matching bounds for estimation error and honest interval length throughout the stated error-scale range. It also compares the neighboring formulas across multiplicative neighborhoods of both thresholds, making the rate characterization uniform through the transitions.

\begin{theoremv}{thm:uniform-frontier}[Uniform risk and length frontier]*
Fix public class constants \(K=(p_{\min},c_{\mathrm{lo}},c_{\mathrm{hi}})\), smoothness exponent \(\beta\), design-decay exponent \(\kappa\), and noncoverage probability \(\alpha\), satisfying:
\begin{itemize}
\item \textbf{(Class constants.)} \(0<p_{\min}\le 1/2\) and \(0<c_{\mathrm{lo}}\le c_{\mathrm{hi}}\).
\item \textbf{(Parameter ranges.)} \(0<\beta\le1\), \(0\le\kappa\le2\), and \(0<\alpha<1/2\).
\item \textbf{(Envelope calibration.)} \(c_{\mathrm{lo}}\le(\kappa+1)2^\kappa\le c_{\mathrm{hi}}\).
\end{itemize}
Write \(R_n(\sigma)\) for the minimax absolute risk in \cref{obj:def:minimax-risk} and \(H_n(\sigma)\) for the minimum honest expected connected-interval length in \cref{obj:def:honest-length}, at these fixed parameters. Define the effective exponent \(d\), logarithmic sample scale \(L_n\), transition thresholds \(a_n,b_n\), and comparison scales \(D_n,F_n(\sigma),P_n(\sigma)\) by
\[
d=2\beta+\kappa+1,\qquad L_n=\log(en),\qquad
a_n=D_n=n^{-1/d},\qquad b_n=L_n^{-1/2},
\]
\[
F_n(\sigma)=\frac{\sigma}{\sqrt{\log(e+n\sigma^d)}},
\qquad
P_n(\sigma)=\frac{\log(e+\sigma^2L_n)}{L_n}.
\]
Define the localization scale \(b_{n,\sigma}\) and absolute-risk rate \(\rho_{n,\sigma}\) by
\[
b_{n,\sigma}=
\begin{cases}
D_n,&\sigma\le a_n,\\
F_n(\sigma),&a_n<\sigma\le b_n,\\
P_n(\sigma),&\sigma>a_n\ \text{and}\ \sigma>b_n,
\end{cases}
\qquad
\rho_{n,\sigma}=b_{n,\sigma}^{\beta}.
\]

There exist constants \(0<c<C\) and an integer \(n_0\), depending on \(K,\beta,\kappa,\alpha\), such that, for every public potential-path space as in \cref{obj:def:potential-path-space}, every integer \(n\ge n_0\), and every \(\sigma\in[0,1/4]\),
\[
c\rho_{n,\sigma}\le R_n(\sigma)\le C\rho_{n,\sigma},
\qquad
c\rho_{n,\sigma}\le H_n(\sigma)\le C\rho_{n,\sigma}.
\]
For the same \(C\) and every \(n\ge n_0\), the transition comparisons satisfy
\[
D_n\le C F_n(a_n),\qquad F_n(a_n)\le C D_n,
\]
\[
F_n(b_n)\le C P_n(b_n),\qquad P_n(b_n)\le C F_n(b_n).
\]

Moreover, for every real \(W\ge1\), there exist a constant \(C_K\ge1\) and an integer \(n_K\) such that, for every integer \(n\ge n_K\),
\[
0<a_n/W,\qquad Wa_n\le1/4,\qquad
0<b_n/W,\qquad Wb_n\le1/4,
\]
and the two transition windows satisfy
\[
C_K^{-1}D_n\le F_n(\sigma)\le C_KD_n
\qquad
\text{for every }\sigma\in[a_n/W,Wa_n],
\]
\[
C_K^{-1}P_n(\sigma)\le F_n(\sigma)\le C_KP_n(\sigma)
\qquad
\text{for every }\sigma\in[b_n/W,Wb_n].
\]
Here \(C_K\) and \(n_K\) may depend on \(W\) and the fixed parameters.

Finally, for every fixed \(\sigma\in(0,1/4]\), there exist constants \(0<c<C\) and an integer \(n_0\), which may depend on \(\sigma\) and the fixed parameters, such that, for every public potential-path space as in \cref{obj:def:potential-path-space} and every integer \(n\ge n_0\),
\[
c\left(\frac{\log\log n}{\log n}\right)^\beta
\le R_n(\sigma)\le
C\left(\frac{\log\log n}{\log n}\right)^\beta,
\]
\[
c\left(\frac{\log\log n}{\log n}\right)^\beta
\le H_n(\sigma)\le
C\left(\frac{\log\log n}{\log n}\right)^\beta.
\]
\end{theoremv}
% CAUSALSMITH-CITED-SCOPE-BEGIN thm:uniform-frontier
\verificationfootnotetext{\textbf{Formalization scope.} This result uses the cited conclusion \citep{PetrushevXu2005} (Theorem 2.4, equation (2.14)), \citep{https://people.math.sc.edu/pencho/Publications/kpx-06-28-06-web.pdf} (Theorem 2.2, equation (2.4)), \citep{https://dlmf.nist.gov/18.3} (Table 18.3.1, Jacobi row and its caption; orthogonality and squared norm.) and \citep{https://dlmf.nist.gov/5.11.E12} (Equation 5.11.12, restricted to the positive real axis.). The cited source proof is not formalized here: Lean verifies its use and all remaining steps. If that cited conclusion is false, the portions of this result that depend on it are not certified.}
% CAUSALSMITH-CITED-SCOPE-END thm:uniform-frontier

The envelope calibration in \cref{obj:thm:uniform-frontier} has a concrete role: the density \((\kappa+1)2^\kappa|t|^\kappa\) integrates to one on \([-1/2,1/2]\), and bracketing its coefficient places the explicit lower-bound design inside the class. Thus the condition certifies a nonempty class containing the witnesses used for the converse.

Under the envelope calibration \(c_{\mathrm{lo}}\le(\kappa+1)2^\kappa\le c_{\mathrm{hi}}\), the three branches express the interaction between local dose support and measurement precision. In the direct regime, the measurement-error scale is at most the localization scale available from an error-free dose record. Local design decay enters through \(d=2\beta+\kappa+1\), and the resulting rate is \(n^{-\beta/d}\). In the intermediate Gaussian regime, recovering a localized mean from the contaminated dose requires Gaussian inversion, with attainable localization scale \(F_n(\sigma)\). For larger error scales, known compact latent support enables polynomial localization: inverse polynomial degree is the localization width, Gaussian inversion increases the inverse weight's second moment with degree, and the largest usable degree has order \(L_n/\log(e+\sigma^2L_n)\). Balancing these effects gives the polynomial localization scale \(P_n(\sigma)\) and explains the transition from the Fourier branch. The detailed public certificate appears in \cref{obj:lem:dictionary-score-rate}. For the calibrated class, raising each localization scale to the smoothness exponent gives both the estimation rate and the honest-length rate.

\Cref{tab:precision-regimes} displays these branches and their thresholds. The scales \(D_n\), \(F_n(\sigma)\), and \(P_n(\sigma)\) are those defined in \cref{obj:thm:uniform-frontier}; each row gives the rate up to the uniform constants in that result.

\begin{table}[htbp]
\centering
\caption{Measurement precision and optimal pointwise uncertainty}
\label{tab:precision-regimes}
\begin{tabular}{lll}
\hline
Regime & Measurement-error range & Risk and honest-length rate \\
\hline
Direct
& \(0\le \sigma\le a_n=n^{-1/d}\)
& \(D_n^\beta=n^{-\beta/d}\) \\[4pt]
Intermediate Gaussian
& \(a_n<\sigma\le b_n=L_n^{-1/2}\)
& \(\displaystyle F_n(\sigma)^\beta
=\left\{\frac{\sigma}{\sqrt{\log(e+n\sigma^d)}}\right\}^{\beta}\) \\[8pt]
Compact support
& \(b_n<\sigma\le1/4\)
& \(\displaystyle P_n(\sigma)^\beta
=\left\{\frac{\log(e+\sigma^2L_n)}{L_n}\right\}^{\beta}\) \\[6pt]
\hline
\end{tabular}
\end{table}

The transition comparisons give the thresholds a precise interpretation. For each fixed multiplicative factor \(W\ge1\), the direct and intermediate scales are comparable throughout \([a_n/W,Wa_n]\) once the sample size is sufficiently large. Likewise, the intermediate and compact-support scales are comparable throughout \([b_n/W,Wb_n]\). Thus neighboring descriptions yield the same order of uncertainty throughout each transition window. The comparison constants may depend on \(W\) and the fixed class parameters, while remaining uniform over the public potential-outcome path spaces.

For every fixed positive measurement-error scale in the stated range, increasing sample size eventually places the experiment in the compact-support regime. Both minimax error and minimum honest expected length then have order \((\log\log n/\log n)^\beta\), with comparison constants that may depend on the fixed error scale. The full uniform characterization also describes sequences of improving measurement precision: attaining the direct-regime order requires an error scale at most of order \(n^{-1/d}\), including its fixed multiplicative transition neighborhood.

The error-free endpoint follows directly from this characterization. The next result states that endpoint for both criteria and places its estimation rate alongside a direct regression benchmark with the same local design-decay exponent.

\begin{propositionv}{prop:error-free-reduction}[Error-free rates and benchmark comparison]*
Fix the public class constants \(K=(p_{\min},c_{\mathrm{lo}},c_{\mathrm{hi}})\) and parameters \(\beta,\kappa,\alpha\) satisfying:
\begin{itemize}
\item \textbf{(Class constants.)} \(0<p_{\min}\le 1/2\) and \(0<c_{\mathrm{lo}}\le c_{\mathrm{hi}}\).
\item \textbf{(Parameter ranges.)} \(0<\beta\le1\), \(0\le\kappa\le2\), and \(0<\alpha<1/2\), where \(\alpha\) is the noncoverage probability.
\item \textbf{(Envelope calibration.)}
\[
c_{\mathrm{lo}}\le(\kappa+1)2^\kappa\le c_{\mathrm{hi}}.
\]
\end{itemize}
There exist constants \(0<c<C\) and \(n_0\in\mathbb N\) such that, for every public potential-path space \((\mathcal S,\mathcal E)\) satisfying \cref{obj:def:potential-path-space} and every integer \(n\ge n_0\), the minimax absolute risk and minimum honest expected connected-interval length of \cref{obj:def:minimax-risk,obj:def:honest-length} satisfy
\[
c\,n^{-\beta/(2\beta+\kappa+1)}
\le R_n(0)\le
C\,n^{-\beta/(2\beta+\kappa+1)},
\qquad
c\,n^{-\beta/(2\beta+\kappa+1)}
\le H_n(0)\le
C\,n^{-\beta/(2\beta+\kappa+1)}.
\]

Moreover, fix any separate direct Gaussian-response benchmark with:
\begin{itemize}
\item \textbf{(Benchmark constants.)} \(c_{\mathrm G},r_{\mathrm G},M_{\mathrm G},\tau_{\mathrm G}>0\).
\item \textbf{(Evaluation neighborhood.)} \(x_*\in(0,1)\) and \(0<\delta\le\min\{x_*,1-x_*\}\).
\item \textbf{(Design.)} A fixed probability distribution \(\nu\) on \([0,1]\) with a nonnegative density \(g\) satisfying
\[
g(u)=c_{\mathrm G}|u-x_*|^\kappa
\quad\text{for }u\in[0,1]\text{ with }|u-x_*|\le\delta.
\]
\end{itemize}
The benchmark observations are \(B_i=f(D_i)+\xi_i\), \(i=1,\ldots,n\), where the \(D_i\) are independent draws from \(\nu\), and the \(\xi_i\) are independent \(N(0,\tau_{\mathrm G}^2)\) errors independent of the design. In the supremum below, \(f\) ranges over Borel functions on \([0,1]\) satisfying
\[
\|f\|_\infty\le M_{\mathrm G},
\qquad
\inf_{\substack{p\text{ a real polynomial}\\ \deg p\le\lfloor\beta\rfloor}}
\ \sup_{\substack{u\in[0,1]\\ |u-x_*|\le a}}
|f(u)-p(u-x_*)|
\le r_{\mathrm G}a^\beta
\]
for every
\[
0<a\le
\left\{\frac{\tau_{\mathrm G}^2(\kappa+1)}
{2c_{\mathrm G}r_{\mathrm G}^2n}\right\}^{1/(2\beta+\kappa+1)}.
\]
The infimum below ranges over measurable estimators from the benchmark observations, and \(E_f\) denotes expectation under their observed law. There exist constants \(0<c<C\) and \(n_0\in\mathbb N\), now also depending on the fixed benchmark, such that for every public potential-path space satisfying \cref{obj:def:potential-path-space} and every integer \(n\ge n_0\),
\[
c\,\inf_{\widehat f}\sup_f E_f|\widehat f(x_*)-f(x_*)|
\le R_n(0)\le
C\,\inf_{\widehat f}\sup_f E_f|\widehat f(x_*)-f(x_*)|.
\]
\end{propositionv}
% CAUSALSMITH-CITED-SCOPE-BEGIN prop:error-free-reduction
\verificationfootnotetext{\textbf{Formalization scope.} This result uses the cited conclusion \citep{PetrushevXu2005} (Theorem 2.4, equation (2.14)), \citep{https://people.math.sc.edu/pencho/Publications/kpx-06-28-06-web.pdf} (Theorem 2.2, equation (2.4)), \citep{https://dlmf.nist.gov/18.3} (Table 18.3.1, Jacobi row and its caption; orthogonality and squared norm.), \citep{https://dlmf.nist.gov/5.11.E12} (Equation 5.11.12, restricted to the positive real axis.) and \citep{Gaiffas2005} (Definition 2 and Theorem 1, equations (2.1), (2.4), and (2.5), PDF pages 3--4; exact power-law specialization with loss parameter one.). The cited source proof is not formalized here: Lean verifies its use and all remaining steps. If that cited conclusion is false, the portions of this result that depend on it are not certified.}
% CAUSALSMITH-CITED-SCOPE-END prop:error-free-reduction

\Cref{obj:prop:error-free-reduction} identifies the precision available when the latent dose is observed exactly. Its exponent reflects both mean smoothness and the declining frequency of observations near the evaluation dose. The comparison with Gaiffas's benchmark is a comparison of minimax absolute-error rates for the displayed function class and fixed benchmark parameters \citep[Definition 2 and Theorem 1, equations (2.1), (2.4), and (2.5), pp.~3--4]{Gaiffas2005}. At \(\beta=1\), that benchmark retains approximation by polynomials of degree at most one. The honest-length conclusion for the causal experiment is supplied by the first part of \cref{obj:prop:error-free-reduction}. Together, these conclusions provide an error-free reference for interpreting the measurement-error branches.

The upper and lower bounds have complementary statistical roles. Observable localization produces estimates of the stratum-specific causal means from contaminated-dose and outcome records; the construction and its honest intervals are given in \cref{obj:thm:observable-upper,obj:thm:finite-honesty}. The converse compares structural laws whose causal means are separated while their induced observed laws remain difficult to distinguish, as stated in \cref{obj:thm:observed-lower}. Working with the observed experiment makes the lower bound apply to every estimator and honest connected interval in the criteria above.

The dependency roadmap is short. Identification in \cref{obj:prop:identification} fixes the observed-law target. The public certificate in \cref{obj:lem:public-error-certificate} turns each inverse weight into a risk and interval-radius bound; the dictionary comparisons in \cref{obj:lem:dictionary-score-rate} select entries at the three frontier scales. Finally, the observed-law witnesses in \cref{obj:thm:observed-lower} convert packet cancellation into matching testing lower bounds. Detailed algebra and witness coordinates are deferred to the two technical appendices.

The weighted-packet construction supporting that converse is developed in \cref{obj:thm:weighted-packet-construction}. Its mathematical ingredients include filtered Jacobi localization \citep[Theorem 2.4, equation (2.14)]{PetrushevXu2005}, the corresponding local difference estimate \citep[Theorem 2.2, equation (2.4)]{https://people.math.sc.edu/pencho/Publications/kpx-06-28-06-web.pdf}, Jacobi orthogonality and squared norms \citep[Table 18.3.1, Jacobi row and caption]{https://dlmf.nist.gov/18.3}, and the positive-real-axis gamma-ratio asymptotic \citep[equation (5.11.12)]{https://dlmf.nist.gov/5.11.E12}. These ingredients supply localization and moment cancellation for the observed-law comparison; their detailed construction belongs with the auxiliary results in the appendix.
